\documentclass[11pt,a4paper]{article}
\usepackage[margin=25mm,headheight=14pt]{geometry}
\usepackage[T1]{fontenc}
\usepackage[utf8]{inputenc}
\usepackage{lmodern,amsmath,amssymb}
\usepackage{xcolor,graphicx,booktabs,longtable,array}
\usepackage{microtype,enumitem,fancyhdr,titlesec}
\usepackage{hyperref}
\definecolor{accent}{HTML}{163B55}
\hypersetup{colorlinks=true,linkcolor=accent,urlcolor=accent,
 pdftitle={A variational spectral lower bound from concentration of separated sets},pdfauthor={Matthew J. Colbrook},
 pdfsubject={AIM 1013: mathematical result and proof}}
\urlstyle{same}
\Urlmuskip=0mu plus 1mu
\setlength{\parindent}{0pt}
\setlength{\parskip}{0.35em}
\setlength{\emergencystretch}{3em}
\setlist{itemsep=0.25em,topsep=0.3em}
\providecommand{\tightlist}{\setlength{\itemsep}{0pt}\setlength{\parskip}{0pt}}
\titleformat{\section}{\normalfont\large\bfseries\color{accent}}{\thesection}{0.7em}{}
\titleformat{\subsection}{\normalfont\normalsize\bfseries}{\thesubsection}{0.6em}{}
\pagestyle{fancy}
\fancyhf{}
\fancyhead[L]{\small\sffamily AIM 1013}
\fancyhead[R]{\small\sffamily Matthew J. Colbrook}
\fancyfoot[C]{\thepage}
\renewcommand{\headrulewidth}{0.3pt}
\setcounter{secnumdepth}{0}
\allowdisplaybreaks[2]
\title{\textbf{A variational spectral lower bound from concentration of separated sets}}
\author{Matthew J. Colbrook\\[3pt]
{\small Department of Applied Mathematics and Theoretical Physics}\\
{\small University of Cambridge, United Kingdom}\\
{\small\href{mailto:mjc249@maths.cam.ac.uk}{\texttt{mjc249@maths.cam.ac.uk}}}}
\date{3 November 2026}
\begin{document}
\maketitle
\begin{abstract}
We derive a quantitative variational spectral lower bound from the multiset concentration hypothesis in AIM 1013, under nonnegative Bakry--\'Emery curvature. The proof uses a half-volume isoperimetric region and separated distance layers, and applies without assuming discrete spectrum.
\end{abstract}
\noindent\textbf{Repository target.}
This manuscript concerns AIM problem 1013 at revision
\texttt{3fa50884d129}. Its scope is stated below.
\medskip

\hypertarget{1013-result}{%
\section{Result}\label{1013-result}}

The answer to Problem 1013 is \textbf{yes}. In fact its concentration
hypothesis implies a lower bound for the first variational gap, and
hence for every higher variational level:

\[
\boxed{
\lambda^{(k)}\geq\lambda^{(1)}
\geq \frac{a}{64k^2(k-1)^2\log 2},\qquad k\geq2.
}
\tag{1}
\]

Here \(a\) is exactly the parameter in the question. The constant is not
claimed to be optimal. The proof applies to complete, possibly
noncompact, smooth connected weighted Riemannian manifolds and does not
require discrete spectrum.

The geometric input is the classical isoperimetric comparison theory
under \(\operatorname{Ric}+\operatorname{Hess}V\geq0\). The additional
step is to divide one side of a half-volume isoperimetric region into
alternating pieces and gaps of prescribed measures. The other half of
the manifold serves as one of the \(k\) pieces.

\hypertarget{1013-isoperimetric-facts-and-their-precise-scope}{%
\section{1. Isoperimetric facts and their precise
scope}\label{1013-isoperimetric-facts-and-their-precise-scope}}

Write

\[
P_\mu(E)=\int_{\partial^*E}e^{-V}\,d\mathcal H^{n-1},
\qquad
\mathcal I(v)=\inf_{\mu(E)=v}P_\mu(E),
\qquad
I=\mathcal I(1/2),
\]

where \(\partial^*E\) is the reduced boundary and \(n=\dim M\). Under
the assumptions of the question, the following facts apply.

\begin{enumerate}
\def\labelenumi{\arabic{enumi}.}
\tightlist
\item
  Isoperimetric minimizers of every measure \(0<v<1\) exist. A minimizer
  can be represented by an open set. Its boundary is smooth outside a
  singular set of Hausdorff dimension at most \(n-8\), and its regular
  part has constant weighted mean curvature. Its weighted perimeter is
  the weighted area of that regular part.
\item
  The profile \(\mathcal I\) is symmetric about \(1/2\) and concave on
  \((0,1)\). Consequently its Cheeger constant is \[
  h_\mu:=\inf_{0<\mu(E)<1}
  \frac{P_\mu(E)}{\min(\mu(E),1-\mu(E))}
  =2\mathcal I(1/2)=2I.
  \tag{2}
  \]
\item
  Normal rays from the regular boundary of an isoperimetric region cover
  its interior up to a set of measure zero. Weighted normal Jacobians
  obey the Heintze--Karcher comparison estimate under the Bakry--Émery
  curvature bound.
\end{enumerate}

These facts hold for the complete finite-measure setting, not just for
compact manifolds. In particular, finite total measure prevents loss of
a prescribed positive volume to infinity in the existence argument:
local BV compactness, followed by the uniform bound
\(\mu(M\setminus K)\to0\) along a compact exhaustion, gives global
convergence in \(L^1(\mu)\). Lower semicontinuity of weighted perimeter
then gives a minimizer. Boundary regularity and the normal-ray assertion
are the standard geometric measure theory and comparison inputs cited
below.

For (2), concavity and nonnegativity imply

\[
\mathcal I(v)\geq2v\mathcal I(1/2)
\quad(0<v\leq1/2).
\]

Symmetry gives the corresponding inequality for \(v\geq1/2\), while a
half-volume minimizer gives equality in the infimum defining \(h_\mu\).

The inputs can be found together in Milman (2010), Sections 2--3:
existence and regularity in Section 2.2, the normal-ray treatment of
singular boundaries in Section 2.4, and Theorem 3.2 and Corollary 3.3
for the profile and Cheeger constant. Milman (2015), Theorem 2.3,
Corollary 2.5, and Proposition 2.9 give a detailed treatment of the
geometric inputs, including the fact that nearest boundary points are
regular. The required comparison is a specialization of its Theorem 2.1.

These statements do not require an orientation of the ambient manifold.
Weighted perimeter uses the Riemannian density; an isoperimetric region
supplies an outward normal on its regular boundary, even if the ambient
manifold is nonorientable. The compactness, regularity, first-variation,
and normal-Jacobian arguments are local and apply in coordinate charts
without an ambient orientation.

The one-dimensional cases are also included. A complete connected
one-dimensional manifold without boundary is a line or a circle. On the
line the density is log-concave, and half-lines are isoperimetric
minimizers; on a circle, \(V''\geq0\) forces \(V\) to be constant and
arcs are minimizers. Their distance-density comparisons below follow
directly from log-concavity, or from the same normal-ray calculation
with no tangential second fundamental form. The Cheeger-profile identity
(2) holds in these cases as well.

\hypertarget{1013-a-half-volume-region-with-bounded-inward-distance-density}{%
\section{2. A half-volume region with bounded inward distance
density}\label{1013-a-half-volume-region-with-bounded-inward-distance-density}}

Choose an open isoperimetric region \(E\) with

\[
\mu(E)=\frac12,
\qquad P_\mu(E)=I.
\]

Its complement, represented by its interior, is another half-volume
minimizer with the opposite weighted mean curvature. Thus we may choose
the side so that the weighted mean curvature with respect to its outward
normal is nonnegative:

\[
H_\mu(E)=\operatorname{div}_{\partial E}\nu-\langle\nabla V,\nu\rangle=H\geq0.
\]

Here the curvature is the trace of the second fundamental form, without
division by \(n-1\), minus the outward normal derivative of \(V\).
Define the globally 1-Lipschitz function

\[
f(x)=d(x,M\setminus E).
\tag{3}
\]

It vanishes on the complement and is positive in \(E\). We claim that,
for all \(0\leq s<t<\infty\),

\[
\mu\{s<f<t\}\leq I(t-s).
\tag{4}
\]

For clarity, (4) is a bound on every distance layer, stronger than a
bound only on the total layer adjacent to the boundary.

To prove it, let \(S\) be the regular part of \(\partial E\), and follow
the inward unit normal from \(x\in S\) until its first cut time
\(\tau(x)\). Let \(J_x(t)\) be the normal Jacobian including the density
ratio, normalized by \(J_x(0)=1\). Before the cut time, the Riccati
equation gives

\[
\frac{d^2}{dt^2}\log J_x(t)
=-|B_t|^2
-\bigl(\operatorname{Ric}+\operatorname{Hess}V\bigr)(\dot\gamma,\dot\gamma)
\leq0.
\]

The initial logarithmic derivative is \(-H\). Hence

\[
J_x(t)\leq e^{-Ht}\leq1,
\qquad 0\leq t<\tau(x).
\tag{5}
\]

The normal-ray representation, including the standard treatment of the
singular boundary and cut locus, now yields

\[
\begin{aligned}
\mu\{s<f<t\}
&\leq\int_S\int_s^t
\mathbf1_{\{r<\tau(x)\}}J_x(r)\,dr\,
 e^{-V(x)}d\mathcal H^{n-1}(x)\\
&\leq(t-s)P_\mu(E)=I(t-s).
\end{aligned}
\]

This proves (4). Equivalently, the pushforward of \(\mu|_E\) under \(f\)
is absolutely continuous on \((0,\infty)\), with density at most \(I\),
and total mass \(1/2\). The quantity \(I\) is finite and positive: a set
of zero weighted perimeter on a connected manifold with positive smooth
density is constant almost everywhere, which is incompatible with
measure \(1/2\).

In dimension one, choose the half-line for which the density is
nonincreasing into its interior. Convexity of \(V\) guarantees one such
side at the median; the density at the median is \(I\), proving (4). For
a circle with uniform measure, the two inward endpoint rays give the
constant density \(I\) up to their meeting point.

\hypertarget{1013-the-separated-pieces}{%
\section{3. The separated pieces}\label{1013-the-separated-pieces}}

Fix \(k\geq2\), and set

\[
b=\frac1{2k},
\qquad
g=\frac{b}{k-1}=\frac1{2k(k-1)}.
\tag{6}
\]

Let

\[
G(t)=\mu\{0<f<t\}.
\]

By (4), \(G\) is continuous and \(I\)-Lipschitz, with \(G(0)=0\) and
\(G(t)\to1/2\) as \(t\to\infty\). Choose finite numbers

\[
0=s_1<t_1<s_2<t_2<\cdots<s_{k-1}<t_{k-1}
\]

such that

\[
G(s_j)=(j-1)(g+b),
\qquad
G(t_j)=jg+(j-1)b,
\qquad 1\leq j\leq k-1.
\tag{7}
\]

All target values in (7) are below \(1/2\), so such thresholds exist.
Flat portions of \(G\) cause no problem, and every positive level set of
\(f\) has measure zero by (4).

Define \(k\) Borel sets by

\[
A_1=\{f=0\},
\]

\[
A_{j+1}=\{t_j\leq f\leq s_{j+1}\}
\quad(1\leq j\leq k-2),
\qquad
A_k=\{f\geq t_{k-1}\}.
\tag{8}
\]

For \(k=2\), the middle family is empty. The omitted gaps are

\[
D_j=\{s_j<f<t_j\},
\qquad 1\leq j\leq k-1,
\]

and each has measure \(g\). Therefore

\[
\mu(A_1)=\frac12,
\qquad
\mu(A_i)=b\quad(2\leq i\leq k),
\qquad
\mu(A)=1-b,
\quad A=\bigcup_{i=1}^kA_i.
\tag{9}
\]

In particular, the measure tuple belongs to the set \(\Delta_k\) in the
question: its sum is \(1-b\), and every individual measure is at least
\(b\).

The density bound (4) and the gap masses give

\[
t_j-s_j\geq\frac{g}{I}.
\]

Since \(f\) is 1-Lipschitz, any two consecutive sets in (8) are at
distance at least the intervening gap width; nonconsecutive sets are
farther apart. Thus

\[
\min_{i\ne j}d(A_i,A_j)\geq\frac{g}{I}>0.
\tag{10}
\]

\hypertarget{1013-testing-the-concentration-assumption}{%
\section{4. Testing the concentration
assumption}\label{1013-testing-the-concentration-assumption}}

Set

\[
r=\frac{g}{4I}.
\tag{11}
\]

By (10), this is within the radius range allowed by the hypothesis:

\[
0<r\leq\frac12\min_{i\ne j}d(A_i,A_j).
\]

If a point of a gap \(D_j\) lies in \(A_r\), the 1-Lipschitz property of
\(f\) forces its \(f\)-value to lie within \(r\) of one of that gap's
two endpoints. Each such endpoint layer has measure at most \(Ir\), by
(4). Consequently

\[
\begin{aligned}
\mu(A_r\setminus A)
&\leq2(k-1)Ir\\
&=\frac{(k-1)g}{2}
=\frac b2.
\end{aligned}
\]

Combining this with (9),

\[
1-\mu(A_r)\geq\frac b2.
\tag{12}
\]

The assumed multi-set concentration estimate, applied to the sets (8),
gives

\[
1-\mu(A_r)
\leq b\exp\bigl[-\min(ar^2,\sqrt a\,r)\bigr].
\]

Together with (12), this implies

\[
\min(ar^2,\sqrt a\,r)\leq\log2.
\]

Writing \(x=\sqrt a\,r\), and using \(\log2<1\), the last inequality
forces \(x<1\) and then \(x^2\leq\log2\). Thus

\[
a\leq\frac{\log2}{r^2}
=\frac{16\log2}{g^2}I^2
=64k^2(k-1)^2\log2\,I^2.
\tag{13}
\]

\hypertarget{1013-from-isoperimetry-to-the-variational-levels}{%
\section{5. From isoperimetry to the variational
levels}\label{1013-from-isoperimetry-to-the-variational-levels}}

The weighted Cheeger inequality on a complete manifold gives

\[
\lambda^{(1)}\geq\frac{h_\mu^2}{4}=I^2.
\tag{14}
\]

This is a functional inequality, so it does not require an eigenfunction
or a discrete spectrum. One may verify it directly: for a real
\(H^1(\mu)\) function \(u\), choose a median \(m\). Apply the coarea
form of the Cheeger inequality to \((u-m)_+^2\) and \((u-m)_-^2\), whose
supports have measure at most \(1/2\), and then use Cauchy--Schwarz.
This gives

\[
\int|\nabla u|^2d\mu
\geq\frac{h_\mu^2}{4}\int|u-m|^2d\mu
\geq\frac{h_\mu^2}{4}\int|u-\mu(u)|^2d\mu.
\]

Every two-dimensional subspace in the variational definition contains a
nonzero mean-zero function, proving (14). Monotonicity of the min--max
levels gives \(\lambda^{(k)}\geq\lambda^{(1)}\). Equations (13)--(14)
prove (1).

If one permits a connected zero-dimensional manifold, it is a single
point and the infimum over \((k+1)\)-dimensional subspaces is empty,
hence \(+\infty\); the assertion is then immediate under the usual
variational convention.

\hypertarget{1013-references-and-verification-scope}{%
\section{References and verification
scope}\label{1013-references-and-verification-scope}}

\begin{enumerate}
\def\labelenumi{\arabic{enumi}.}
\tightlist
\item
  N. Gozlan and R. Herry, \emph{Multiple sets exponential concentration
  and higher order eigenvalues}, Potential Analysis \textbf{52} (2020),
  203--221, Section 5.2. DOI:
  \url{https://doi.org/10.1007/s11118-018-9743-1}. Primary preprint:
  \url{https://arxiv.org/abs/1804.06133}.
\item
  E. Milman, \emph{Isoperimetric and Concentration Inequalities ---
  Equivalence under Curvature Lower Bound}, Duke Mathematical Journal
  \textbf{154} (2010), 207--239. Sections 2--3, especially Section 2.2,
  Section 2.4, Theorem 3.2, and Corollary 3.3. DOI:
  \url{https://doi.org/10.1215/00127094-2010-038}. Primary final
  preprint: \url{https://arxiv.org/pdf/0902.1560}.
\item
  E. Milman, \emph{Sharp isoperimetric inequalities and model spaces for
  the Curvature-Dimension-Diameter condition}, Journal of the European
  Mathematical Society \textbf{17} (2015), 1041--1078. Theorem 2.1,
  Theorem 2.3, Corollary 2.5, and Proposition 2.9. DOI:
  \url{https://doi.org/10.4171/JEMS/526}. Primary published text:
  \url{https://ems.press/content/serial-article-files/32036}.
\item
  E. Milman, \emph{Uniform tail-decay of Lipschitz functions implies
  Cheeger's isoperimetric inequality under convexity assumptions},
  Comptes Rendus Mathématique \textbf{346} (2008), 989--994, Theorem 4.2
  and the ensuing Cheeger-profile identity. DOI:
  \url{https://doi.org/10.1016/j.crma.2008.07.022}. Primary published
  text:
  \url{https://www.numdam.org/item/10.1016/j.crma.2008.07.022.pdf}.
\end{enumerate}
\end{document}
